\section[Методы и модели классификации при дисбалансе классов]{Методы и модели классификации при дисбалансе классов в прикладных задачах машинного обучения}
\label{sec:chapter1}

\subsection{Постановка задачи бинарной и многоклассовой классификации при перекосе частот}
\label{sec:ch1-1}

\subsection{Метрики качества и критерии сравнения моделей при дисбалансе классов}
\label{sec:ch1-2}

\subsection{Методы уровня данных: перевыборка, синтез примеров и фильтрация}
\label{sec:ch1-3}

\subsection{Методы уровня алгоритма: взвешивание, cost-sensitive learning и специализированные функции потерь}
\label{sec:ch1-4}

\subsection{Ансамблевые, каскадные и двухэтапные архитектуры}
\label{sec:ch1-5}

\subsection{Особенности данных и постановок задач в прикладных доменах}
\label{sec:ch1-6}

\subsection{Сравнительный анализ существующих решений и нерешённые проблемы}
\label{sec:ch1-7}

\subsection{Требования к предметно-ориентированному подходу и постановка задач исследования}
\label{sec:ch1-8}
